\chapter{Beliefs, Rationality and Strategic Choice}
\label{sec:p17-epistemic-game-theory}

A choice can be optimal for a player's beliefs even when those beliefs
are mistaken. Epistemic game theory asks which choices follow from
specified beliefs about other players, their beliefs, and their
reasoning. The finite-type presentation developed here follows Perea's
chapter on epistemic game theory
\cite[chap.~9.2, pp.~543--551]{Perea2021EpistemicGameTheory}.
It gives the complete finite static construction, its elimination
procedure, both numerical examples, and the chapter's dynamic example.
The necessary finite proofs are supplied in the annex.

The source contrasts classical solution concepts, which select possible
strategies, with conditions on the beliefs that support a choice.
It presents Harsanyi's types as a compact representation of belief
hierarchies and connects the elimination result to the rationalizability
work of Bernheim and Pearce. These historical connections identify
different presentations; they do not identify common belief with
correct prediction of what other players actually do.

\section{A finite model of beliefs about beliefs}

\begin{hbdefinition}{def:p17-egt-types}{Game, types and subjective beliefs}
Let \(I\) be a finite nonempty player set. Each \(C_i\) is finite
and nonempty, \(C=\prod_i C_i\), and \(u_i:C\to\mathbb R\).
These are the normal-form data of
\cref{def:p17-finite-strategic-game}, with \(C_i\) replacing \(A_i\).
The game and its utility interpretation are fixed throughout the model.
Every type uses these same action sets and payoff functions; the
uncertainty concerns choices and beliefs, rather than the identity
of the game or its utilities.

A finite epistemic model consists of nonempty finite sets \(T_i\)
and probability distributions
\[
 b_i(t_i)\in\Delta\!\left(\prod_{j\ne i}(C_j\times T_j)\right)
 \qquad(t_i\in T_i).
\]
A type specifies beliefs; it is not an action. Write \(q_i(t_i)\)
for the marginal distribution on \(C_{-i}=\prod_{j\ne i}C_j\).
This joint distribution need not factor into independent opponents'
distributions. Different players' beliefs need not come from one
common prior, and a positive-probability conjecture need not be true.
An empty product, for a one-player game, is a singleton.

Define the expected payoff and optimal choices by
\[
 v_i(c_i,t_i)=\sum_{c_{-i}}q_i(t_i)(c_{-i})u_i(c_i,c_{-i}),
 \qquad
 B_i(t_i)=\mathop{\rm argmax}_{d_i\in C_i}v_i(d_i,t_i).
\]
Every \(B_i(t_i)\) is nonempty. A realized choice is rational for its
type when it belongs to this set. Maximizing a linear expected utility
does not require randomization: some pure choice is optimal, although
mixtures of tied optima are also optimal if the decision domain permits them.
\index{epistemic game theory}\index{type!epistemic game}
\end{hbdefinition}

\begin{hbtheorem}{thm:p17-egt-hierarchies}{proved=annex:p17.egt.hierarchies}{A finite type model determines every finite belief level}
Set \(h_i^0(t_i)=*\). Recursively, \(h_i^{r+1}(t_i)\) is the
distribution of
\[
                  (c_j,h_j^r(t_j))_{j\ne i}
\]
under \(b_i(t_i)\). Each is a well-defined finite-support probability
distribution, and forgetting the last belief level recovers the
preceding one. Thus each type determines a coherent infinite sequence
of finite-level beliefs without writing an infinite input table.
This constructs the hierarchies of this model; it does not assert
that one finite model contains every possible hierarchy.
\end{hbtheorem}

\section{Belief in rationality and common belief}

\begin{hbdefinition}{def:p17-egt-common-belief}{Individual levels and cumulative levels}
Let \(H_i^1\) consist of types \(t_i\) for which every
positive-probability tuple \((c_j,t_j)_{j\ne i}\) satisfies
\(c_j\in B_j(t_j)\) for all \(j\ne i\).
For \(r\geq1\), define
\[
 H_i^{r+1}=\{t_i:b_i(t_i)\text{ gives probability one to tuples
                         with }t_j\in H_j^r\text{ for every }j\ne i\}.
\]
These are the source's onefold, twofold and higher individual levels.
For reasoning that requires \emph{all} levels through \(r\), put
\[
 Q_i^0=T_i,\qquad Q_i^r=\bigcap_{\ell=1}^{r}H_i^\ell,\qquad
 Q_i^\infty=\bigcap_{r\geq1}Q_i^r .
\]
A type expresses common belief in rationality when it belongs to
\(Q_i^\infty\). A choice is rationally possible under that condition
when it is optimal for such a type in some finite epistemic model.
This existential condition differs from asking whether it is optimal
at a particular given type.
\end{hbdefinition}

\begin{hbtheorem}{thm:p17-egt-belief-algorithm}{proved=annex:p17.egt.belief.algorithm}{Finite graph computation and the distinction between levels}
Make a directed graph on the tagged types \((i,t_i)\), with an edge
to \((j,t_j)\) whenever that opponent type occurs with positive
marginal probability under \(b_i(t_i)\). Mark a type bad if it is
outside \(H_i^1\). A type has common belief in rationality exactly
when no bad type is reachable from it, counting the starting vertex.
Equivalently,
\[
 Q_i^{r+1}=H_i^1\cap
       \{t_i:\text{every successor belongs to }Q_j^r\}.
\]
This decreasing computation stabilizes after at most
\(\sum_i|T_i|\) strict removal rounds. Individual levels \(H_i^r\)
need not be decreasing: twofold belief alone does not imply onefold
belief. A complete two-choice example is given in the proof.
\end{hbtheorem}

\section{The elimination theorem}

\begin{hbdefinition}{def:p17-egt-reduction}{Strict domination in a restricted game}
For nonempty subsets \(R_i\subseteq C_i\), write \(R=\prod_iR_i\).
A choice \(a_i\in R_i\) is strictly dominated in \(R\) if some
\(x_i\in\Delta(R_i)\) satisfies
\[
 \sum_{d_i\in R_i}x_i(d_i)u_i(d_i,c_{-i})
                 >u_i(a_i,c_{-i})
                 \quad\text{for every }c_{-i}\in R_{-i}.
\]
The inequality is strict at every remaining opponent profile.
The mixture is a comparison device; no random action is required
of the player whose choices are being characterized.
Set \(R_i^0=C_i\) and simultaneously retain in \(R_i^{k+1}\)
exactly the choices not strictly dominated in \(R^k\).
Put \(R_i^\infty=\bigcap_kR_i^k\).
\index{rationalizability}\index{strict domination!iterated elimination}
\end{hbdefinition}

\begin{hblemma}{thm:p17-egt-best-response-alternative}{proved=annex:p17.egt.best.response.alternative}{Belief witness or strict domination witness}
In any such finite restriction, a choice \(a_i\in R_i\) is
undominated exactly when there is a probability
\(p_i\in\Delta(R_{-i})\) for which it maximizes expected payoff
over all choices in \(R_i\).
More explicitly, define
\[
 D_{d_i,c_{-i}}=u_i(d_i,c_{-i})-u_i(a_i,c_{-i}).
\]
Exactly one of the following holds:
\[
\begin{aligned}
 &\exists x_i\in\Delta(R_i)\quad
       \min_{c_{-i}}\sum_{d_i}x_i(d_i)D_{d_i,c_{-i}}>0;\\
 &\exists p_i\in\Delta(R_{-i})\quad Dp_i\leq0 .
\end{aligned}
\]
The second witness may correlate different opponents' choices.
The proof uses the fully established finite minimax theorem
\cref{thm:p17-finite-minimax}. This is the alternative attributed
to Pearce's Lemma 3 in the source chapter.
\end{hblemma}

\begin{hbtheorem}{thm:p17-egt-reduction-integrity}{proved=annex:p17.egt.reduction.integrity}{Nonempty restrictions and the original deviation set}
Every \(R_i^k\) is nonempty, and the simultaneous procedure stabilizes
after at most \(\sum_i(|C_i|-1)\) strict removal rounds.
For every belief supported on \(R_{-i}^k\), all its best responses
in the original \(C_i\) belong to \(R_i^k\).
Consequently a surviving action's restricted best-response witness
is also a witness against \emph{all original deviations}.
No deleted action is silently ignored in the epistemic optimality test.
\end{hbtheorem}

\begin{hbtheorem}{thm:p17-egt-finite-depth}{proved=annex:p17.egt.finite.depth}{The exact number of reasoning levels}
For \(k\geq0\) and \(a_i\in C_i\), the following are equivalent:
\begin{enumerate}
\item \(a_i\in R_i^{k+1}\).
\item In some finite epistemic model, \(a_i\in B_i(t_i)\)
      for a type \(t_i\in Q_i^k\).
\end{enumerate}
Thus rational choice without an additional belief requirement
corresponds to the first elimination round; rational choice plus
onefold belief corresponds to two rounds. The construction in the
proof supplies finite types and all their beliefs explicitly.
\end{hbtheorem}

\begin{hbtheorem}{thm:p17-egt-fundamental}{proved=annex:p17.egt.fundamental}{Common belief in rationality and iterated strict elimination}
A choice is rationally possible under common belief in rationality
if and only if it belongs to \(R_i^\infty\).
For the reverse direction one can use exactly one type \(t_i^a\)
per remaining choice \(a\in R_i^\infty\), attaching a supporting
belief to each choice and pairing every conjectured opponent
choice with its corresponding type. All types in the resulting
finite model express common belief in rationality.
\end{hbtheorem}

\begin{hbmethod}{def:p17-egt-computation}{Computing the answer and its witnesses}
Supply the finite action lists and cardinal payoff arrays. For each
current action solve the displayed domination or belief inequalities.
Delete all strictly dominated actions simultaneously and repeat to
stability. Return the remaining action sets, every deletion's mixture
and positive margin, and a supporting probability for each final action.
These probabilities build the finite epistemic witness model.

For rational payoff data this is a finite sequence of rational linear
programs: maximize a margin \(z\) subject to
\(x_i\geq0\), \(\sum x_i=1\), and \(x_i^{\mathsf T}D\geq z\mathbf1\).
Remove an action precisely when the optimum is positive. The dual
best-response inequalities provide a retention witness.
For an explicit terminating procedure, bound \(0\leq z\leq M\),
where \(M=\max|D_{d_i,c_{-i}}|\), enumerate full-rank sets of
active constraints, solve them by rational elimination, and retain
the feasible solutions with largest \(z\).
When that largest value is zero, enumerate the vertices of
\(\{p_i\geq0:\sum p_i=1,\ Dp_i\leq0\}\) to obtain a belief.
The proof of \cref{thm:p17-egt-best-response-alternative} supplies
the existence of these vertices and the exact enumeration argument.
Arbitrary real inputs require an input representation and an exact
comparison procedure before this is an effective algorithm.
No psychological prediction follows merely from executing the method.
\end{hbmethod}

\section{The source's elimination example}

\begin{hbdefinition}{def:p17-egt-first-game}{Payoffs and six types}
The source's Table 9.2.1 has the following payoff pairs:
\[
\begin{array}{c|rrrr}
 &e&f&g&h\\\hline
a&(0,0)&(4,1)&(4,4)&(4,3)\\
b&(3,2)&(0,0)&(3,4)&(3,3)\\
c&(2,2)&(2,1)&(0,0)&(2,3)\\
d&(1,2)&(1,1)&(1,4)&(0,0)
\end{array}
\]
The types are \(t_1^a,t_1^b,t_1^c\) and \(t_2^e,t_2^g,t_2^h\).
A written pair below denotes its point mass:
\[
\begin{array}{ll}
b_1(t_1^a)=(g,t_2^g),&
b_2(t_2^e)=\tfrac35(c,t_1^c)+\tfrac25(d,t_1^a),\\
b_1(t_1^b)=(e,t_2^e),&
b_2(t_2^g)=(a,t_1^a),\\
b_1(t_1^c)=\tfrac35(e,t_2^e)+\tfrac25(f,t_2^g),&
b_2(t_2^h)=(c,t_1^c).
\end{array}
\]
Superscripts label the optimal choice for each type in this example.
They do not add a general axiom identifying types and choices.
\end{hbdefinition}

\begin{hbtheorem}{thm:p17-egt-first-game}{proved=annex:p17.egt.first.game}{Every payoff, belief level and deletion in the example}
The six displayed superscripts are their types' unique optimal choices.
Exactly \(t_1^a,t_1^b,t_2^g,t_2^h\) express onefold belief;
only \(t_1^a,t_2^g\) express all levels through two and common belief.
The elimination sequence is
\[
 \begin{array}{c|cc}
 k&R_1^k&R_2^k\\\hline
 0&\{a,b,c,d\}&\{e,f,g,h\}\\
 1&\{a,b,c\}&\{e,g,h\}\\
 2&\{a,b\}&\{g,h\}\\
 3&\{a\}&\{g\}.
 \end{array}
\]
The annex computes every type's expected-payoff vector, gives each
deletion witness, and verifies that the retained choices have witnesses
at the round in which they are retained.
\end{hbtheorem}

\section{What Nash equilibrium adds}

\begin{hbtheorem}{thm:p17-egt-nash-beliefs}{proved=annex:p17.egt.nash.beliefs}{The hierarchy induced by an equilibrium}
Let \(\sigma\) be a Nash equilibrium of the finite game, as in
\cref{def:p17-nash-equilibrium}. Give each player one type \(t_i\)
and set
\[
 b_i(t_i)((c_j,t_j)_{j\ne i})=\prod_{j\ne i}\sigma_j(c_j).
\]
Every type expresses common belief in rationality, and every action
in the support of \(\sigma_i\) is optimal for \(t_i\).
Here the equilibrium probabilities may be interpreted as conjectures
about choices. This interpretation alone does not show that actual
players randomize, share a prior, or correctly predict one another.
\end{hbtheorem}

\begin{hbdefinition}{def:p17-egt-gap-game}{A game with a unique equilibrium and more rationalizable choices}
The source's Table 9.2.3 is
\[
 \begin{array}{c|rrr}
 &d&e&f\\\hline
 a&(0,3)&(3,0)&(0,2)\\
 b&(3,0)&(0,3)&(0,2)\\
 c&(2,0)&(2,0)&(2,2).
 \end{array}
\]
Use six types with the deterministic beliefs
\[
\begin{aligned}
 t_1^a&\mapsto(e,t_2^e),& t_2^e&\mapsto(b,t_1^b),\\
 t_1^b&\mapsto(d,t_2^d),& t_2^d&\mapsto(a,t_1^a),\\
 t_1^c&\mapsto(f,t_2^f),& t_2^f&\mapsto(c,t_1^c).
\end{aligned}
\]
These arrows specify belief assignments, including the conjectured choice.
\end{hbdefinition}

\begin{hbtheorem}{thm:p17-egt-nash-gap}{proved=annex:p17.egt.nash.gap}{Complete equilibrium calculation and the epistemic gap}
All six types in this model express common belief in rationality,
and each superscript is its type's unique optimal choice.
Thus every pure choice is rationally possible under common belief.
The only Nash equilibrium, including mixed profiles, is \((c,f)\).
At \(t_1^a\), player 1 believes that player 2 chooses \(e\), but
believes that player 2 believes that player 1 believes that player 2
chooses \(d\). Common belief in rationality therefore permits this
specific error about another player's first-order belief.
\end{hbtheorem}

\begin{hbtheorem}{thm:p17-egt-mutual-conjectures}{proved=annex:p17.egt.mutual.conjectures}{A precise two-player condition yielding Nash conjectures}
Fix actual types \(t_1,t_2\) in a two-player model. Assume both
believe in the other's rationality. Also assume that each assigns
probability one to opponent types having the opponent's
\emph{actual first-order belief}:
\[
 b_i(t_i)(c_j,s_j)>0\quad\Longrightarrow\quad
                 q_j(s_j)=q_j(t_j)\qquad(i\ne j).
\]
Then \((q_2(t_2),q_1(t_1))\) is a Nash equilibrium.
This statement concerns the two actual types together, with the
correctness condition explicitly quantified. It does not replace
correctness by agreement about the payoff matrix or by additional
unqualified iterations of rationality alone.
\end{hbtheorem}

The original Aumann--Brandenburger conditions are supplied in
Chapter~\ref{sec:p17-epistemic-nash-conditions}.
Its \cref{thm:p17-ab-two-player} treats mutual knowledge of the
game, rationality and full conjectures for two players.
Its \cref{thm:p17-ab-many-player} treats a common prior assigning
positive probability to mutual knowledge of the game and rationality
together with common knowledge of the full conjectures, yielding
agreement, independence and Nash for many players.
The exact finite fixed-game bridge in both directions is
\cref{thm:p17-ab-type-bridge}. That treatment preserves the
distinction between equilibrium conjectures, actual actions
and null states. It does not identify these source conditions
with the separate Polak or one-person perspective/depth
presentations \cite{AumannBrandenburger1995Epistemic}.

\begin{hbtheorem}{thm:p17-egt-many-player-boundaries}{proved=annex:p17.egt.many.player.boundaries}{Correlation and disagreement are independent additional issues}
With at least three players, common belief in rationality need
not make a player's joint conjecture about opponents independent.
Even independent conjectures need not make two players agree
about a third player's action distribution.
The annex gives complete finite-type countermodels for both claims.
In the first, replacing a correlated belief by the product of its
marginals reverses the player's unique optimal choice.
\end{hbtheorem}

\section{Belief revision in a dynamic game}

\begin{hbdefinition}{def:p17-egt-reny-game}{The complete four-move example}
In the source's Figure 9.2.2, player 1 chooses \(a\) or \(b\);
\(a\) ends with payoff \((3,3)\). After \(b\), player 2 chooses
\(c\) or \(d\); \(c\) ends with \((2,2)\). After \(bd\), player 1
chooses \(e\) or \(f\); \(e\) ends with \((1,1)\).
After \(bdf\), player 2 chooses \(g\) or \(h\), ending with
\((0,4)\) or \((4,0)\), respectively. The history is observed at
every decision, and payoffs are ordered by player.

A pure strategy specifies a choice at both of a player's nodes:
\(C_1=\{ae,af,be,bf\}\) and \(C_2=\{cg,ch,dg,dh\}\).
A terminal outcome is distinct from this full contingent strategy.
\end{hbdefinition}

\begin{hbtheorem}{thm:p17-egt-reny-reasoning}{proved=annex:p17.egt.reny.reasoning}{Backward reasoning and a different response to the surprising move}
Backward optimization chooses \(g,e,c,a\), in that reverse order,
and produces contingent strategies \((ae,cg)\) and outcome \(a\).

Suppose instead that after observing \(b\), player 2 insists,
whenever possible, on explaining it by an originally optimal
pure strategy of player 1 for some belief about player 2's
strategies. The only such strategy compatible with \(b\) is \(bf\).
Expecting its continuation \(f\), player 2 chooses \(d\), then \(g\)
if called upon again. This is a different off-path plan from
backward reasoning. If player 1 anticipates that plan and
optimizes, the initial outcome is still \(a\).

To make \(bf\) originally optimal, a conjecture with probabilities
\(p_c,p_{dg},p_{dh}\) on stopping at \(c\), choosing \(dg\), and
choosing \(dh\) must satisfy
\[
 p_c+p_{dg}+p_{dh}=1,\qquad
                 2p_c+4p_{dh}\geq3 .
\]
In particular, conditional on continuation at \(d\), it assigns
probability at least \(3/4\) to \(h\).
The claim is about the explicitly stated explanatory rule at this
history; it is not a general definition of extensive-form rationalizability.
\end{hbtheorem}

\begin{proofboundary}{def:p17-egt-source-boundaries}{Representations and dynamic refinements with distinct obligations}
The static theorem uses the source's full finite game and finite-type
domain, allowing arbitrary joint conjectures. Existence of terminal
type structures containing all hierarchies is a separate result;
the finite construction does not establish it. Likewise, a
state-based encoding needs its own precise interpretation and
translation, not just renaming states as types.

The source compares common belief in future rationality, which permits
past mistakes while retaining specified future-rationality beliefs,
with strong belief in rationality, which retains rational explanations
of observed play whenever compatible ones remain. It relates the first
to backward and sequential reasoning and the second to forward
induction, explicable equilibrium and extensive-form rationalizability.
It also reports a general coincidence of backward and forward
\emph{outcomes} for perfect-information games without relevant ties.
These cited general claims require the corresponding conditional-belief
models and definitions of relevant ties; the calculation above proves
the displayed game only. The single-player reformulations of the
two-player conjecture conditions and their exact finite belief-depth
requirements also retain their source-specific definitions.
The original chapter is the identified source of this comparison
\cite[pp.~545, 548--550]{Perea2021EpistemicGameTheory};
the full underlying source presentations remain separate work.
\end{proofboundary}
